% !TEX root = ../main.tex
\chapter{Pretrained Weights and Hub}
\label{ch:pretrained}

A model zoo's value is in large part the pretrained weights it
ships --- a randomly initialised ResNet-50 is just a tensor
shape, but with ImageNet weights it is a usable feature
extractor. \Lucid\ ships pretrained checkpoints for every
model in the zoo through its \emph{weights hub}: a small,
content-addressed download protocol that handles version
management, format conversion, sharded downloads, and offline
caching. This chapter documents the hub's design, the on-disk
format, the download protocol, and the
\code{from\_pretrained} contract that the model factories
honour.

\section{Goals}
\label{sec:pretrained:goals}

The hub design targets three goals:

\begin{itemize}
  \item \textbf{Reproducibility}: the same hub call returns
    the same weights, indefinitely. New weights get new
    version tags; old weights are not silently overwritten.
  \item \textbf{Fast first call}: a typical
    \code{from\_pretrained("resnet\_50")} should fetch and load
    in under 10 seconds on a fast network.
  \item \textbf{Robustness}: partial downloads resume; corrupt
    files are detected by hash and re-fetched.
\end{itemize}

The hub does not target:
\begin{itemize}
  \item Hosting arbitrary user-uploaded weights (\Lucid\
    delegates this to HuggingFace Hub, with a thin
    \code{lucid.hub.from\_huggingface} bridge).
  \item Versioning user-trained weights (the user owns
    \code{state\_dict.save}/\code{load}; the hub is for
    \Lucid-curated checkpoints).
\end{itemize}

\section{Hub URL Structure}
\label{sec:pretrained:url_structure}

Each pretrained checkpoint has a canonical URL:

\begin{lstlisting}[style=lucid,language=LucidPython]
https://hub.lucid-ml.org/{family}/{model_id}/{tag}/weights.safetensors
\end{lstlisting}

Components:
\begin{itemize}
  \item \code{family}: e.g.\ \code{vision/cnn},
    \code{language/bert}.
  \item \code{model\_id}: the canonical model name, e.g.\
    \code{resnet\_50}, \code{bert\_base}.
  \item \code{tag}: pretraining dataset + version, e.g.\
    \code{imagenet1k-v1}, \code{coco-v2}.
\end{itemize}

The full URL is appended with the file format suffix
(\code{.safetensors} is the default; \code{.npz} legacy is
supported as fallback).

\section{The safetensors Format}
\label{sec:pretrained:safetensors}

\Lucid\ uses safetensors \cite{safetensors} as the primary
on-disk format. The format consists of:

\begin{itemize}
  \item An 8-byte length-prefixed JSON header listing tensors
    by name with their dtype, shape, and byte offsets.
  \item The raw tensor data, packed contiguously.
\end{itemize}

The format's advantages over the pickle-based
\cite{pickle_python} formats of legacy zoos:

\begin{enumerate}
  \item \textbf{Memory-mapped loading}: the file can be mmap'd
    and tensors materialised on-demand. \Lucid's tensor
    factory wraps this directly: a 7\,GB LLaMA checkpoint
    loads in milliseconds (header parse only).
  \item \textbf{Safe deserialisation}: no arbitrary code
    execution. Pickle-based formats allow loading malicious
    weights that execute on read.
  \item \textbf{Cross-framework compatibility}: the format is
    a published spec; checkpoints can be read by any compatible
    library.
\end{enumerate}

\subsection{Loader API}
\label{ssec:pretrained:safetensors_loader}

\begin{lstlisting}[style=lucid,language=LucidPython]
from lucid.serialization import load_safetensors

state_dict = load_safetensors("model.safetensors")
# state_dict is a dict[str, Tensor]; tensors are mmap'd lazily

model.load_state_dict(state_dict)
\end{lstlisting}

The lazy materialisation matters: a partial load (e.g., loading
only the encoder of a Mask-R-CNN model) reads only the requested
tensors' bytes from disk, saving most of the I/O.

\section{The Download Protocol}
\label{sec:pretrained:download_protocol}

\code{from\_pretrained} downloads via a small in-house client
(no external HTTP libraries are required beyond Python's stdlib
\code{urllib}):

\begin{enumerate}
  \item Compute the cache path:
    \code{\$LUCID\_HOME/hub/<family>/<model\_id>/<tag>/weights.safetensors}.
    (\code{LUCID\_HOME} defaults to \code{\$HOME/.lucid}.)
  \item If the file exists and its SHA-256 matches the
    expected hash (fetched from a tiny manifest), use it.
  \item Otherwise, download with HTTP Range requests in
    8\,MiB chunks; on failure, resume from the last
    complete chunk.
  \item After download, verify the SHA-256 again. If it
    fails, delete and re-download (one retry); on second
    failure, raise.
\end{enumerate}

\subsection{Hash manifest}
\label{ssec:pretrained:hash_manifest}

Each model has a manifest JSON:

\begin{lstlisting}[style=lucid,language=LucidPython]
{
    "model_id": "resnet_50",
    "tag": "imagenet1k-v1",
    "sha256": "a3b4c5d6...",
    "size_bytes": 102457288,
    "format": "safetensors",
    "files": [
        {"name": "weights.safetensors", "sha256": "...", "size": ...},
    ],
}
\end{lstlisting}

The manifest is fetched separately (~1\,KB) before any large
download. It includes:
\begin{itemize}
  \item The expected SHA-256 (for hash check).
  \item File listing (for multi-file checkpoints like sharded
    LLaMA).
  \item Format hint (\code{safetensors} or \code{npz}).
\end{itemize}

\section{Sharded Checkpoints}
\label{sec:pretrained:sharded}

For models above 4\,GB, the checkpoint is split into multiple
files (typically 1\,GB each):

\begin{itemize}
  \item \code{weights-00001-of-00007.safetensors}
  \item \code{weights-00002-of-00007.safetensors}
  \item \dots
  \item \code{index.json} mapping parameter names to file shard
    + offset.
\end{itemize}

The loader reads \code{index.json}, opens all shards mmap-style,
and presents a unified \code{state\_dict} interface. Shard
boundaries do not respect parameter boundaries; one large
parameter may span multiple shards (with the per-shard segment
correctly assembled by the loader).

\subsection{Parallel download}
\label{ssec:pretrained:parallel_download}

Shards are downloaded in parallel (up to 4 concurrent
connections). For LLaMA-7B (13\,GB across 7 shards), this
roughly quarters the wall-clock vs sequential download on a
high-bandwidth connection.

\section{AutoModel.from\_pretrained}
\label{sec:pretrained:auto_model}

The user-facing API:

\begin{lstlisting}[style=lucid,language=LucidPython]
model = lucid.models.AutoModel.from_pretrained("resnet_50")
\end{lstlisting}

Internally, the call dispatches:
\begin{enumerate}
  \item Looks up the model in the registry
    (\chref{ch:zoo_architecture}~Section~\ref{sec:zoo_architecture:register_model}).
    Errors if unknown.
  \item Looks up the default \code{pretrained\_on}
    (\code{imagenet1k-v1} for ResNet-50).
  \item Fetches the manifest, downloads the weights if not
    cached, loads the safetensors.
  \item Instantiates the model with the default config.
  \item Calls \code{model.load\_state\_dict(weights, strict=True)}.
\end{enumerate}

Override the dataset:
\begin{lstlisting}[style=lucid,language=LucidPython]
model = lucid.models.AutoModel.from_pretrained(
    "vit_base_16", pretrained_on="imagenet21k",
)
\end{lstlisting}

The set of valid \code{pretrained\_on} values per model is
listed in the registry; an invalid value errors with the
allowed list.

\subsection{Strict vs non-strict loading}
\label{ssec:pretrained:strict_load}

By default, \code{from\_pretrained} loads strict: every weight
in the checkpoint must match a parameter in the model and vice
versa. For fine-tuning with a modified head:

\begin{lstlisting}[style=lucid,language=LucidPython]
model = lucid.models.AutoModel.from_pretrained(
    "resnet_50",
    strict=False,            # Allow missing/extra keys
    ignore_mismatched_sizes=True,  # Allow tensor-size mismatches
)
\end{lstlisting}

The non-strict load reports the keys that did not match (so the
user knows what was re-initialised) but does not raise.

\section{Fine-tuning Workflow}
\label{sec:pretrained:finetune_workflow}

Typical fine-tuning starts with a pretrained backbone and a
fresh head:

\begin{lstlisting}[style=lucid,language=LucidPython]
import lucid.models as lm
import lucid.nn as nn
import lucid.optim as optim

# Load backbone with frozen weights
model = lm.AutoModel.from_pretrained("resnet_50")
for p in model.parameters():
    p.requires_grad = False

# Replace classifier head
model.fc = nn.Linear(2048, num_classes)  # ImageNet to your dataset

# Train only the new head
optimizer = optim.Adam(
    [p for p in model.parameters() if p.requires_grad],
    lr=1e-3,
)

for x, y in dataloader:
    out = model(x)
    loss = nn.functional.cross_entropy(out, y)
    loss.backward()
    optimizer.step()
\end{lstlisting}

For full-model fine-tuning (don't freeze), use a smaller
learning rate (typically $10^{-5}$) to avoid catastrophic
forgetting.

\section{Offline Cache}
\label{sec:pretrained:offline_cache}

The hub respects the standard environment variables for offline
use:

\begin{itemize}
  \item \code{LUCID\_OFFLINE=1}: skip all network calls; only
    use cached weights. Error if not cached.
  \item \code{LUCID\_HOME}: cache directory (default
    \code{\$HOME/.lucid}).
  \item \code{LUCID\_HUB\_URL}: override the hub URL (for
    self-hosted mirrors).
\end{itemize}

Pre-download for offline use:

\begin{lstlisting}[style=lucid,language=LucidPython]
lucid.hub.prefetch([
    "resnet_50",
    "bert_base",
    ("vit_base_16", "imagenet21k"),
])
# Downloads and caches all without instantiating models
\end{lstlisting}

\section{HuggingFace Bridge}
\label{sec:pretrained:hf_bridge}

For weights not in the \Lucid\ hub, the
\code{from\_huggingface} bridge fetches via the HuggingFace API
and converts the format:

\begin{lstlisting}[style=lucid,language=LucidPython]
model = lucid.models.bert.from_huggingface(
    "bert-base-uncased",
    revision="main",
)
\end{lstlisting}

The bridge:
\begin{enumerate}
  \item Downloads \code{config.json}, \code{pytorch\_model.bin}
    (or safetensors variant).
  \item Reads the config, instantiates the matching \Lucid\
    model.
  \item Translates the weight names from the HuggingFace
    convention to \Lucid's (per-family translation table).
  \item Loads.
\end{enumerate}

Translation tables live in
\code{lucid/models/<family>/\_hf\_keymap.py} as a simple
\code{dict[str, str]}.

\section{Format Conversion Tools}
\label{sec:pretrained:conversion}

\Lucid\ ships a command-line tool for the inverse conversion:

\begin{lstlisting}[style=lucid,language=LucidPython]
python -m lucid.tools.export \
    --model resnet_50 \
    --weights ./my_finetuned.safetensors \
    --format onnx \
    --output ./model.onnx
\end{lstlisting}

Supported export formats: ONNX \cite{onnx_2017} (for
cross-framework deployment), Core ML (for Apple-platform
deployment), safetensors (for re-uploading to \Lucid\ hub).

The ONNX exporter traces a forward pass and emits the equivalent
ONNX graph. Caveat: not all \Lucid\ ops have ONNX equivalents
(custom Functions in particular); the exporter raises clearly
on unsupported ops.

\section{Version Pinning}
\label{sec:pretrained:version_pinning}

Each pretrained-weights tag is immutable. Improving the
checkpoint of a model (e.g., a longer training run) gets a new
tag:

\begin{itemize}
  \item \code{imagenet1k-v1}: original release weights.
  \item \code{imagenet1k-v2}: refreshed weights (better
    augmentation, longer training).
\end{itemize}

Default tags point to the recommended version:

\begin{lstlisting}[style=lucid,language=LucidPython]
# In lucid/models/vision/cnn/resnet/_pretrained.py:
@register_model(
    ...
    default_pretrained="imagenet1k-v2",  # Latest stable
)
def resnet_50(config): ...
\end{lstlisting}

A user explicitly pinning a version:

\begin{lstlisting}[style=lucid,language=LucidPython]
model = lucid.models.AutoModel.from_pretrained(
    "resnet_50", pretrained_on="imagenet1k-v1",
)
\end{lstlisting}

Old versions remain downloadable indefinitely; the hub does
not garbage-collect.

\section{Hub Statistics}
\label{sec:pretrained:hub_stats}

The current \Lucid\ hub holds:

\begin{table}[t]
\centering
\small
\begin{tabular}{lrr}
\toprule
Family & Models & Total weights (GB) \\
\midrule
Vision CNN          & 38 & 11.2 \\
Vision Transformers & 18 & 14.5 \\
Detection / Seg     & 21 & 18.7 \\
Language --- BERT   & 14 & 4.8  \\
Language --- GPT2   & 12 & 6.1  \\
Language --- T5     & 8  & 12.3 \\
Generative          & 11 & 25.4 \\
\midrule
Total               & 122 & 93.0 \\
\bottomrule
\end{tabular}
\caption[Lucid hub statistics.]{Number of pretrained checkpoints
hosted in the \Lucid\ hub, by family. Counts include multiple
\code{pretrained\_on} variants per model. Storage is on a CDN
fronted by R2; per-file requests are typically served in
\(<\)100\,ms first-byte from the US west coast.}
\label{tab:pretrained:hub_stats}
\end{table}

The numbers grow each release as new families and as v2/v3
weight refreshes accumulate.

\section{Failure Modes}
\label{sec:pretrained:failures}

Common failure modes and the hub's handling:

\begin{itemize}
  \item \textbf{Network failure}: retries with exponential
    backoff, capped at 3 attempts. After 3 failures, raises
    \code{HubDownloadError} with the cause.
  \item \textbf{Disk full}: download fails on write; partial
    file is deleted; \code{HubCacheError} is raised with the
    free-space report.
  \item \textbf{Hash mismatch}: file deleted and one re-download
    attempted; second mismatch raises
    \code{HubIntegrityError} (the user's connection is
    suspect).
  \item \textbf{Stale registry}: model name not in registry
    but in hub. The error suggests
    \code{lucid.hub.list\_available()} for the current
    registry.
\end{itemize}

\section{Summary and Forward References}
\label{sec:pretrained:summary}

The pretrained-weights hub provides reproducible, fast,
robust access to \Lucid's curated checkpoints. The safetensors
format gives fast lazy loading; sharded downloads cover the
multi-gigabyte language models; the HuggingFace bridge fills
the gap for non-curated weights.

This chapter closes \partref{part:model_zoo}.
\partref{part:performance} treats the performance-engineering
work that all of these models depend on for usable training
throughput: methodology, AMP integration, MPSGraph dispatch
tuning, and the no-compile and compile speedup sweeps.

\begin{lucidnote}
For the user: \code{from\_pretrained} is the canonical entry
point; for offline use, set \code{LUCID\_OFFLINE=1} and
\code{lucid.hub.prefetch} the needed weights ahead of time.
For weights from elsewhere, prefer the safetensors format and
\code{load\_safetensors} directly.
\end{lucidnote}
